% Options for packages loaded elsewhere
% Options for packages loaded elsewhere
\PassOptionsToPackage{unicode}{hyperref}
\PassOptionsToPackage{hyphens}{url}
\PassOptionsToPackage{dvipsnames,svgnames,x11names}{xcolor}
%
\documentclass[
  letterpaper,
  DIV=11,
  numbers=noendperiod]{scrartcl}
\usepackage{xcolor}
\usepackage{amsmath,amssymb}
\setcounter{secnumdepth}{5}
\usepackage{iftex}
\ifPDFTeX
  \usepackage[T1]{fontenc}
  \usepackage[utf8]{inputenc}
  \usepackage{textcomp} % provide euro and other symbols
\else % if luatex or xetex
  \usepackage{unicode-math} % this also loads fontspec
  \defaultfontfeatures{Scale=MatchLowercase}
  \defaultfontfeatures[\rmfamily]{Ligatures=TeX,Scale=1}
\fi
\usepackage{lmodern}
\ifPDFTeX\else
  % xetex/luatex font selection
\fi
% Use upquote if available, for straight quotes in verbatim environments
\IfFileExists{upquote.sty}{\usepackage{upquote}}{}
\IfFileExists{microtype.sty}{% use microtype if available
  \usepackage[]{microtype}
  \UseMicrotypeSet[protrusion]{basicmath} % disable protrusion for tt fonts
}{}
% Make \paragraph and \subparagraph free-standing
\makeatletter
\ifx\paragraph\undefined\else
  \let\oldparagraph\paragraph
  \renewcommand{\paragraph}{
    \@ifstar
      \xxxParagraphStar
      \xxxParagraphNoStar
  }
  \newcommand{\xxxParagraphStar}[1]{\oldparagraph*{#1}\mbox{}}
  \newcommand{\xxxParagraphNoStar}[1]{\oldparagraph{#1}\mbox{}}
\fi
\ifx\subparagraph\undefined\else
  \let\oldsubparagraph\subparagraph
  \renewcommand{\subparagraph}{
    \@ifstar
      \xxxSubParagraphStar
      \xxxSubParagraphNoStar
  }
  \newcommand{\xxxSubParagraphStar}[1]{\oldsubparagraph*{#1}\mbox{}}
  \newcommand{\xxxSubParagraphNoStar}[1]{\oldsubparagraph{#1}\mbox{}}
\fi
\makeatother


\usepackage{longtable,booktabs,array}
\usepackage{calc} % for calculating minipage widths
% Correct order of tables after \paragraph or \subparagraph
\usepackage{etoolbox}
\makeatletter
\patchcmd\longtable{\par}{\if@noskipsec\mbox{}\fi\par}{}{}
\makeatother
% Allow footnotes in longtable head/foot
\IfFileExists{footnotehyper.sty}{\usepackage{footnotehyper}}{\usepackage{footnote}}
\makesavenoteenv{longtable}
\usepackage{graphicx}
\makeatletter
\newsavebox\pandoc@box
\newcommand*\pandocbounded[1]{% scales image to fit in text height/width
  \sbox\pandoc@box{#1}%
  \Gscale@div\@tempa{\textheight}{\dimexpr\ht\pandoc@box+\dp\pandoc@box\relax}%
  \Gscale@div\@tempb{\linewidth}{\wd\pandoc@box}%
  \ifdim\@tempb\p@<\@tempa\p@\let\@tempa\@tempb\fi% select the smaller of both
  \ifdim\@tempa\p@<\p@\scalebox{\@tempa}{\usebox\pandoc@box}%
  \else\usebox{\pandoc@box}%
  \fi%
}
% Set default figure placement to htbp
\def\fps@figure{htbp}
\makeatother





\setlength{\emergencystretch}{3em} % prevent overfull lines

\providecommand{\tightlist}{%
  \setlength{\itemsep}{0pt}\setlength{\parskip}{0pt}}



 


\usepackage{booktabs}
\usepackage{caption}
\usepackage{longtable}
\usepackage{colortbl}
\usepackage{array}
\usepackage{anyfontsize}
\usepackage{multirow}
\usepackage{xcolor, soul}
\definecolor{Light}{gray}{.90}
\sethlcolor{Light}
\let\OldTexttt\texttt
\renewcommand{\texttt}[1]{\OldTexttt{\hl{#1}}}
\usepackage{caption}
\captionsetup{font=small, labelfont=bf}
\KOMAoption{captions}{tableheading}
\makeatletter
\@ifpackageloaded{caption}{}{\usepackage{caption}}
\AtBeginDocument{%
\ifdefined\contentsname
  \renewcommand*\contentsname{Table of contents}
\else
  \newcommand\contentsname{Table of contents}
\fi
\ifdefined\listfigurename
  \renewcommand*\listfigurename{List of Figures}
\else
  \newcommand\listfigurename{List of Figures}
\fi
\ifdefined\listtablename
  \renewcommand*\listtablename{List of Tables}
\else
  \newcommand\listtablename{List of Tables}
\fi
\ifdefined\figurename
  \renewcommand*\figurename{Figure}
\else
  \newcommand\figurename{Figure}
\fi
\ifdefined\tablename
  \renewcommand*\tablename{Table}
\else
  \newcommand\tablename{Table}
\fi
}
\@ifpackageloaded{float}{}{\usepackage{float}}
\floatstyle{ruled}
\@ifundefined{c@chapter}{\newfloat{codelisting}{h}{lop}}{\newfloat{codelisting}{h}{lop}[chapter]}
\floatname{codelisting}{Listing}
\newcommand*\listoflistings{\listof{codelisting}{List of Listings}}
\makeatother
\makeatletter
\makeatother
\makeatletter
\@ifpackageloaded{caption}{}{\usepackage{caption}}
\@ifpackageloaded{subcaption}{}{\usepackage{subcaption}}
\makeatother
\usepackage{bookmark}
\IfFileExists{xurl.sty}{\usepackage{xurl}}{} % add URL line breaks if available
\urlstyle{same}
\hypersetup{
  pdftitle={Progressive decrease in effective dimensionality of visual cortex regions},
  pdfauthor={Al W. Xin},
  colorlinks=true,
  linkcolor={blue},
  filecolor={Maroon},
  citecolor={Blue},
  urlcolor={Blue},
  pdfcreator={LaTeX via pandoc}}


\title{Progressive decrease in effective dimensionality of visual cortex
regions}
\author{Al W. Xin}
\date{Feb 6, 2026}
\begin{document}
\maketitle


Advisors:

\begin{itemize}
\tightlist
\item
  Statistics: Rob Kass
\item
  External: Jenelle Feather, Assistant Professor in Psychology
  Department and Neuroscience Institute, CMU. Jenelle is my co-advisor
  in the Stats-PNC program.
\end{itemize}

\section{Introduction}\label{introduction}

Creating artificial neural networks (ANNs) that are aligned to observed
brain activity could provide a testbed to perform neuroscience
experiments that are infeasible in living subjects.

Though strategies exist to quantify brain-ANN similarity {[}1{]}, it is
unclear how the underlying structure of data influences these metrics.
Differences in the effective dimensionality of neural responses or model
weights, for example, may influence the current approach to similarity
scoring.

Characterizing the dimensionality of brain regions is contentious, with
investigators drawing contrasting conclusions about whether the visual
cortex is better characterized as high- or low-dimensional {[}2, 3{]}.
The source of divergence remains unclear, and possible causes include
differences in neural recording, task complexity, and choice of
statistical method. Furthermore, the effect of dimensionality on model
similarity is also controversial, with investigators disagreeing about
whether high effective dimensionality drives higher similarity scores
{[}4, 5{]}.

Within this project, I will perform a comprehensive investigation of
data collected in the macaque visual cortex {[}6{]}. Because researchers
presented upwards of 1000 unique stimuli and recorded activity from
different visual regions from the same subjects, I anticipate that this
data set will contain sufficient detail to address existing questions
about neural dimensionality.

\section{Data}\label{data}

Lee et al.~trained 8 macaque monkeys to fixate on a screen and presented
the animals with a passive fixation tasks. The monkeys viewed synthetic
textures samples with natural structure (``naturalistic'') or without
structure (``noise'') {[}7{]}. A subset of monkeys also performed a
4-choice task to discriminate naturalistic and noise samples, though
this behavioral data will likely not be relevant for the project.
Monkeys viewed images in blocks of 8, with images shown for 100 ms and
separated by 100 ms of blank stimulus.

Lee et al.~placed 96-electrode ``Utah'' recording arrays in 8 macaque
monkeys, producing data from visual regions V1, V2, V4, IT, and FC
(foveal confluence). For monkeys M2, M3, and M6, a second electrode
produced parallel recordings in two visual regions, seen in
Table~\ref{tbl-subjects}. Lee et al.~published the binned spiking data,
summed for bins of 10 ms, but omitted the raw spike timing data.

\begin{table}

\caption{\label{tbl-subjects}Description of the macaque monkey subjects
and relevant features of the recording data. For monkeys with more than
one recording over time, the youngest and oldest age recording, in
weeks, is reported as Age 1 and Age 2, respectively.}

\centering{

\fontsize{12.0pt}{14.0pt}\selectfont
\begin{tabular*}{\linewidth}{@{\extracolsep{\fill}}llrrl}
\toprule
Monkey ID & Sex & Age 1 & Age 2 & Region(s) \\ 
\midrule\addlinespace[2.5pt]
M1 & F & 17 & 95 & V1-V2 border, V4 \\ 
M2 & F & 16 & 95 & V2, V4 \\ 
M3 & M & 16 & 65 & FC, IT \\ 
M4 & F & 27 & 53 & NA \\ 
M5 & M & 17 & 64 & NA \\ 
M6 & M & 184 & 194 & V1, IT \\ 
M7 & M & 184 & 194 & IT \\ 
M8 & M & NA & NA & V4 \\ 
\bottomrule
\end{tabular*}

}

\end{table}%

\section{Sources}\label{sources}

\begin{enumerate}
\def\labelenumi{\arabic{enumi}.}
\tightlist
\item
  Schrimpf M, Kubilius J, Lee MJ, Ratan Murty NA, Ajemian R, DiCarlo JJ.
  Integrative Benchmarking to Advance Neurally Mechanistic Models of
  Human Intelligence. Neuron. 2020 Nov 11;108(3):413-423. doi:
  10.1016/j.neuron.2020.07.040. Epub 2020 Sep 11. PMID: 32918861.
\item
  Stringer C, Pachitariu M, Steinmetz N, Carandini M, Harris KD.
  High-dimensional geometry of population responses in visual cortex.
  Nature. 2019 Jul;571(7765):361--5.
\item
  Gao P, Trautmann E, Yu B, Santhanam G, Ryu S, Shenoy K, Ganguli S. A
  theory of multineuronal dimensionality, dynamics and measurement
  {[}Internet{]}. bioRxiv; 2017 {[}cited 2026 Feb 4{]}. p.~214262.
  Available from: https://www.biorxiv.org/content/10.1101/214262v2
\item
  Elmoznino E, Bonner MF. High-performing neural network models of
  visual cortex benefit from high latent dimensionality. PLOS
  Computational Biology. 2024 Jan 10;20(1):e1011792.
\item
  Canatar A, Feather J, Wakhloo A, Chung S. A Spectral Theory of Neural
  Prediction and Alignment. Advances in Neural Information Processing
  Systems. 2023 Dec 15;36:47052--80.
\item
  Lee GM, Deliz CLR, Bushnell BN, Majaj NJ, Movshon JA, Kiorpes L.
  Developmentally stable representations of naturalistic image structure
  in macaque visual cortex. Cell Reports {[}Internet{]}. 2024 Aug 27
  {[}cited 2025 Oct 31{]};43(8). Available from:
  https://www.cell.com/cell-reports/abstract/S2211-1247(24)00863-5
\item
  Portilla J, Simoncelli EP. A Parametric Texture Model Based on Joint
  Statistics of Complex Wavelet Coefficients. International Journal of
  Computer Vision. 2000 Oct 1;40(1):49--70.
\end{enumerate}




\end{document}
